\section{Related Work}
\label{sec:related}

\subsection{Classical Supercompilation and Generalization}
The foundation of supercompilation was laid by Valentin Turchin~\cite{turchin1986concept} in the context of the Refal programming language. Turchin demonstrated that executing functional programs symbolically with positive information propagation could achieve dramatic optimizations, including partial evaluation and loop fusion. However, Refal relied on dynamic Lisp-like string rewriting without a typed memory model or formal termination argument.

Sørensen and Glück~\cite{sorensen1995algorithm} formalized \textit{positive supercompilation} for a first-order functional language, proving that homeomorphic embedding guarantees termination over finite signatures. Klimov~\cite{klimov2008generalization} extended this with systematic generalization strategies, establishing the algebraic properties of Most Specific Generalization (MSG). NumLang incorporates textbook MSG anti-unification but significantly extends it to handle mutable pointer locations, SSA basic block boundaries, and register remapping.

\subsection{Functional Supercompilers (Supero, HOSC, Distillation)}
Bolingbroke and Peyton Jones developed Supero~\cite{bolingbroke2010supercompilation}, an ambitious call-by-need supercompiler for Core Haskell. Supero introduced a well-quasi-ordering termination proof based on term size and tag tracking. However, Supero suffered from severe code size explosion on complex recursive programs and lacked support for native parallel code emission or mechanized verification.

Mitchell and Runciman introduced HOSC~\cite{mitchell2010hosc,klyuchnikov2010hosc}, exploring higher-order supercompilation and proof verification via transformation. Klyuchnikov and Romanenko~\cite{klyuchnikov2012practical} subsequently introduced Multi-Result Supercompilation (MRSC), recognizing that greedy generalization often selects suboptimal residual graphs. NumLang integrates MRSC hypergraph exploration directly into its SSA pipeline, utilizing multi-objective Pareto frontiers to balance code size, loop count, and execution speed.

Hamilton proposed \textit{Distillation}~\cite{hamilton2007distillation}, a global program transformation technique that folds configurations across different function definitions and call sites. Distillation eliminates intermediate data structures that escape local supercompilation (such as nested reversals in \texttt{double\_nrev}). NumLang adapts Hamilton's global fold into its SSA MIR driver, enabling cross-procedural deforestation over imperative control flow.

\subsection{Partial Evaluation, Deforestation, and Stream Fusion}
Supercompilation is intimately related to partial evaluation~\cite{jones1993partial}. While standard partial evaluation evaluates known static expressions and residualizes dynamic ones, supercompilation goes further by exploring speculative branches, accumulating symbolic constraints, and synthesizing new recursive loop topologies. NumLang realizes all three Futamura projections~\cite{futamura1971partial}, achieving self-application and automated compiler generator synthesis on typed SSA intermediate representations.

Wadler's deforestation~\cite{wadler1990deforestation} syntactically eliminates intermediate trees in functional programs constrained to \textit{treeless form}. Gill, Launchbury, and Peyton Jones introduced short-cut deforestation via \texttt{build}/\texttt{foldr}. Coutts et al.~\cite{coutts2007stream} generalized this with \textit{stream fusion}, eliminating intermediate list constructors by transforming recursive generators into coinductive step functions. NumLang achieves stream fusion and deforestation as emergent properties of symbolic driving: by representing closures symbolically and evaluating intermediate consumer loops against producer configurations, intermediate structures are contracted into SSA registers without requiring manual rewrite rules or syntactic restrictions.
